% GENERATED FILE. Edit canonical lessons, then run npm run generate:books.
% canonical-source-hash: fnv1a64:5454e900f6158dd4
% canonical-lessons: TE-C477-rtupavanam, TE-C477-varsapatam, TE-C477-musuru, TE-C477-jallu, TE-C477-tadisipovu, TE-R477-review-470-477

\chapter{The Monsoon}
\label{ch:te-the-monsoon}
\hlchaptermodality{477}

\begin{chapteropening}
\textbf{By the end of this chapter:} \emph{I can say the monsoon, rainfall, a spell of drizzle, a shower of rain and to get soaked.}

Complete the last lesson of chapter 477: I can say the monsoon, rainfall, a spell of drizzle, a shower of rain and to get soaked.
\end{chapteropening}

\section[\texorpdfstring{\te{ఋతుపవనం} (ṛtupavanaṁ)}{ఋతుపవనం (ṛtupavanaṁ)}]{\te{ఋతుపవనం} (ṛtupavanaṁ) — the monsoon}

\label{lesson:TE-C477-rtupavanam}



\begin{quote}
Before the new one: say the Telugu for muggy heat, then the Telugu for to heat up.
\end{quote}

\subsection*{You'll want to know: \te{ఋతుపవనం}}
\textbf{\te{ఋతుపవనం}} — \emph{ṛtupavanaṁ} — \textquotedblleft{}the monsoon\textquotedblright{}.

\begin{quote}
\te{వేసవి} \te{తరువాత} \te{ఋతుపవనాలు} \te{వస్తాయి}.
\end{quote}

\begin{quote}
\emph{vēsavi taruvāta ṛtupavanālu vastāyi.}
\end{quote}

\begin{quote}
After the summer the monsoon arrives.
\end{quote}

Mostly plural: \te{ఋతుపవనాలు}, the monsoon winds. \te{ఋతువు}, a season, + \te{పవనం}, a wind.

\begin{cousinweb}
\textbf{\te{ఋతుపవనం}} (\emph{ṛtupavanaṁ}) is a modern compound built from two Sanskrit words: \emph{ṛtu}, \textquotedblleft{}a season\textquotedblright{}, and \emph{pavana}, \textquotedblleft{}wind, breeze\textquotedblright{} (Monier-Williams): the wind of the season.
\end{cousinweb}

\begin{grammarlens}[title={What you've built}]
One more word for getting around.
\end{grammarlens}

\subsection*{Your turn}
\begin{itemize}
\raggedright
  \item \emph{Say it:} \emph{ṛtupavanaṁ}
  \item \emph{Say it:} \emph{ṛtupavanaṁ}, once more
  \item \emph{Say it:} \emph{vēsavi taruvāta ṛtupavanālu vastāyi.} — After the summer the monsoon arrives.
  \item \emph{Recall:} say the Telugu for muggy heat, then the Telugu for to heat up, then say \emph{ṛtupavanaṁ} again
\end{itemize}

\begin{tcolorbox}[breakable,colback=teal!4,colframe=teal!35!black,title={Before you move on}]
What does \te{ఋతుపవనం} mean? (\textquotedblleft{}The monsoon\textquotedblright{}.) Say it once more.
\end{tcolorbox}
\section[\texorpdfstring{\te{వర్షపాతం} (varṣapātaṁ)}{వర్షపాతం (varṣapātaṁ)}]{\te{వర్షపాతం} (varṣapātaṁ) — rainfall}

\label{lesson:TE-C477-varsapatam}



\begin{quote}
Before the new one: say the Telugu for to heat up, then the Telugu for the monsoon.
\end{quote}

\subsection*{You'll want to know: \te{వర్షపాతం}}
\textbf{\te{వర్షపాతం}} — \emph{varṣapātaṁ} — \textquotedblleft{}rainfall\textquotedblright{}.

\begin{quote}
\te{ఆ} \te{సంవత్సరం} \te{వర్షపాతం} \te{చాలా} \te{తక్కువ}.
\end{quote}

\begin{quote}
\emph{ā saṁvatsaraṁ varṣapātaṁ cālā takkuva.}
\end{quote}

\begin{quote}
Rainfall that year was very low.
\end{quote}

\te{వర్ష}, rain, + \te{పాతం}, a fall: the amount of rain that falls.

\begin{cousinweb}
\textbf{\te{వర్షపాతం}} (\emph{varṣapātaṁ}) is a compound of two words borrowed from Sanskrit: \emph{varṣa}, \textquotedblleft{}rain\textquotedblright{} (CDIAL 11392), and \emph{pāta}, \textquotedblleft{}a fall, a falling\textquotedblright{} (Monier-Williams).
\end{cousinweb}

\begin{grammarlens}[title={What you've built}]
One more word for getting around.
\end{grammarlens}

\subsection*{Your turn}
\begin{itemize}
\raggedright
  \item \emph{Say it:} \emph{varṣapātaṁ}
  \item \emph{Say it:} \emph{varṣapātaṁ}, once more
  \item \emph{Say it:} \emph{ā saṁvatsaraṁ varṣapātaṁ cālā takkuva.} — Rainfall that year was very low.
  \item \emph{Recall:} say the Telugu for to heat up, then the Telugu for the monsoon, then say \emph{varṣapātaṁ} again
\end{itemize}

\begin{tcolorbox}[breakable,colback=teal!4,colframe=teal!35!black,title={Before you move on}]
What does \te{వర్షపాతం} mean? (\textquotedblleft{}Rainfall\textquotedblright{}.) Say it once more.
\end{tcolorbox}
\section[\texorpdfstring{\te{ముసురు} (musuru)}{ముసురు (musuru)}]{\te{ముసురు} (musuru) — a spell of drizzle}

\label{lesson:TE-C477-musuru}



\begin{quote}
Before the new one: say the Telugu for the monsoon, then the Telugu for rainfall.
\end{quote}

\subsection*{You'll want to know: \te{ముసురు}}
\textbf{\te{ముసురు}} — \emph{musuru} — \textquotedblleft{}a spell of drizzle\textquotedblright{}.

\begin{quote}
\te{మూడు} \te{రోజులుగా} \te{ముసురు} \te{పట్టింది}.
\end{quote}

\begin{quote}
\emph{mūḍu rōjulugā musuru paṭṭindi.}
\end{quote}

\begin{quote}
For three days there has been a steady drizzle.
\end{quote}

Grey days of light rain that will not stop: \te{ముసురు} \te{పట్టింది}.

\begin{cousinweb}
\textbf{\te{ముసురు}} (\emph{musuru}) is native Dravidian: DEDR 4914 lists Telugu \emph{musuru}, \textquotedblleft{}a constant or continued rain\textquotedblright{}.
\end{cousinweb}

\begin{grammarlens}[title={What you've built}]
One more word for getting around.
\end{grammarlens}

\subsection*{Your turn}
\begin{itemize}
\raggedright
  \item \emph{Say it:} \emph{musuru}
  \item \emph{Say it:} \emph{musuru}, once more
  \item \emph{Say it:} \emph{mūḍu rōjulugā musuru paṭṭindi.} — For three days there has been a steady drizzle.
  \item \emph{Recall:} say the Telugu for the monsoon, then the Telugu for rainfall, then say \emph{musuru} again
\end{itemize}

\begin{tcolorbox}[breakable,colback=teal!4,colframe=teal!35!black,title={Before you move on}]
What does \te{ముసురు} mean? (\textquotedblleft{}A spell of drizzle\textquotedblright{}.) Say it once more.
\end{tcolorbox}
\section[\texorpdfstring{\te{జల్లు} (jallu)}{జల్లు (jallu)}]{\te{జల్లు} (jallu) — a shower of rain}

\label{lesson:TE-C477-jallu}



\begin{quote}
Before the new one: say the Telugu for rainfall, then the Telugu for a spell of drizzle.
\end{quote}

\subsection*{You'll want to know: \te{జల్లు}}
\textbf{\te{జల్లు}} — \emph{jallu} — \textquotedblleft{}a shower of rain\textquotedblright{}.

\begin{quote}
\te{సాయంత్రం} \te{ఒక} \te{చిన్న} \te{జల్లు} \te{పడింది}.
\end{quote}

\begin{quote}
\emph{sāyantraṁ oka cinna jallu paḍindi.}
\end{quote}

\begin{quote}
In the evening there was a short shower.
\end{quote}

\te{జల్లు} \te{పడు}, a shower falls; \te{చిరుజల్లు} is a light one.

\begin{cousinweb}
\textbf{\te{జల్లు}} (\emph{jallu}) is native Dravidian, of the family of Telugu \emph{callu}, \textquotedblleft{}to sprinkle, to scatter\textquotedblright{} (DEDR 2384).
\end{cousinweb}

\begin{grammarlens}[title={What you've built}]
One more word for getting around.
\end{grammarlens}

\subsection*{Your turn}
\begin{itemize}
\raggedright
  \item \emph{Say it:} \emph{jallu}
  \item \emph{Say it:} \emph{jallu}, once more
  \item \emph{Say it:} \emph{sāyantraṁ oka cinna jallu paḍindi.} — In the evening there was a short shower.
  \item \emph{Recall:} say the Telugu for rainfall, then the Telugu for a spell of drizzle, then say \emph{jallu} again
\end{itemize}

\begin{tcolorbox}[breakable,colback=teal!4,colframe=teal!35!black,title={Before you move on}]
What does \te{జల్లు} mean? (\textquotedblleft{}A shower of rain\textquotedblright{}.) Say it once more.
\end{tcolorbox}
\section[\texorpdfstring{\te{తడిసిపోవు} (taḍisipōvu)}{తడిసిపోవు (taḍisipōvu)}]{\te{తడిసిపోవు} (taḍisipōvu) — to get soaked}

\label{lesson:TE-C477-tadisipovu}



\begin{quote}
Before the new one: say the Telugu for a spell of drizzle, then the Telugu for a shower of rain.
\end{quote}

\subsection*{You'll want to know: \te{తడిసిపోవు}}
\textbf{\te{తడిసిపోవు}} — \emph{taḍisipōvu} — \textquotedblleft{}to get soaked\textquotedblright{}.

\begin{quote}
\te{గొడుగు} \te{లేక} \te{వర్షంలో} \te{తడిసిపోయాను}.
\end{quote}

\begin{quote}
\emph{goḍugu lēka varṣaṁlō taḍisipōyānu.}
\end{quote}

\begin{quote}
I had no umbrella and got soaked in the rain.
\end{quote}

Past (I …): \te{తడిసిపోయాను}, \emph{taḍisipōyānu}.

\te{తడి}, wet, + \te{పోవు}: soaked through. \te{తడిసి} \te{ముద్దయ్యాను}, I got drenched.

\begin{cousinweb}
\textbf{\te{తడిసిపోవు}} (\emph{taḍisipōvu}) is native Dravidian: \te{తడియు}, \textquotedblleft{}to get wet\textquotedblright{}, the verb of \te{తడి}, \textquotedblleft{}wet\textquotedblright{} (DEDR 3045), with \te{పోవు}, \textquotedblleft{}to go\textquotedblright{} (DEDR 4572).
\end{cousinweb}

\begin{grammarlens}[title={What you've built}]
One more word for getting around.
\end{grammarlens}

\subsection*{Your turn}
\begin{itemize}
\raggedright
  \item \emph{Say it:} \emph{taḍisipōvu}
  \item \emph{Say it:} \emph{taḍisipōvu}, once more
  \item \emph{Say it:} \emph{goḍugu lēka varṣaṁlō taḍisipōyānu.} — I had no umbrella and got soaked in the rain.
  \item \emph{Recall:} say the Telugu for a spell of drizzle, then the Telugu for a shower of rain, then say \emph{taḍisipōvu} again
\end{itemize}

\begin{tcolorbox}[breakable,colback=teal!4,colframe=teal!35!black,title={Before you move on}]
What does \te{తడిసిపోవు} mean? (\textquotedblleft{}To get soaked\textquotedblright{}.) Say it once more.
\end{tcolorbox}
\section[(dialogue)]{Words from chapters 470–477 — a review}

\label{lesson:TE-R477-review-470-477}



\begin{quote}
Say the Telugu for pollution and waste.
\end{quote}

\subsection*{Your turn}
Say these aloud:

\begin{itemize}
\raggedright
  \item pollution, waste, a gas, to pollute, pure
  \item the environment, greenery, to chop down, to damage, to conserve
  \item a wild animal, a wildlife sanctuary, to hunt, to die out, to migrate
  \item an isle, a whale, coral, a tidal surge, to roar
  \item a peak, a plateau, a continent, a volcano, peaceful
  \item a confluence, a dam, a reservoir, irrigation, a marsh
  \item the climate, temperature, humidity, muggy heat, to heat up
  \item the monsoon, rainfall, a spell of drizzle, a shower of rain, to get soaked
\end{itemize}

\begin{tcolorbox}[breakable,colback=teal!4,colframe=teal!35!black,title={Before you move on}]
Pollution? (\textbf{\te{కాలుష్యం}}, \emph{kāluṣyaṁ}.) Waste? (\textbf{\te{వ్యర్థాలు}}, \emph{vyarthālu}.) A gas? (\textbf{\te{వాయువు}}, \emph{vāyuvu}.) To pollute? (\textbf{\te{కలుషితం} \te{చేయు}}, \emph{kaluṣitaṁ cēyu}.) Pure? (\textbf{\te{స్వచ్ఛమైన}}, \emph{svacchamaina}.) The environment? (\textbf{\te{పర్యావరణం}}, \emph{paryāvaraṇaṁ}.) Greenery? (\textbf{\te{పచ్చదనం}}, \emph{paccadanaṁ}.) To chop down? (\textbf{\te{నరుకు}}, \emph{naruku}.) To damage? (\textbf{\te{నష్టపరచు}}, \emph{naṣṭaparacu}.) To conserve? (\textbf{\te{సంరక్షించు}}, \emph{saṁrakṣiñcu}.) A wild animal? (\textbf{\te{వన్యప్రాణి}}, \emph{vanyaprāṇi}.) A wildlife sanctuary? (\textbf{\te{అభయారణ్యం}}, \emph{abhayāraṇyaṁ}.) To hunt? (\textbf{\te{వేటాడు}}, \emph{vēṭāḍu}.) To die out? (\textbf{\te{అంతరించిపోవు}}, \emph{antariñcipōvu}.) To migrate? (\textbf{\te{వలసపోవు}}, \emph{valasapōvu}.) An isle? (\textbf{\te{ద్వీపం}}, \emph{dvīpaṁ}.) A whale? (\textbf{\te{తిమింగలం}}, \emph{timiṅgalaṁ}.) Coral? (\textbf{\te{పగడం}}, \emph{pagaḍaṁ}.) A tidal surge? (\textbf{\te{ఉప్పెన}}, \emph{uppena}.) To roar? (\textbf{\te{గర్జించు}}, \emph{garjiñcu}.) A peak? (\textbf{\te{శిఖరం}}, \emph{śikharaṁ}.) A plateau? (\textbf{\te{పీఠభూమి}}, \emph{pīṭhabhūmi}.) A continent? (\textbf{\te{ఖండం}}, \emph{khaṇḍaṁ}.) A volcano? (\textbf{\te{అగ్నిపర్వతం}}, \emph{agniparvataṁ}.) Peaceful? (\textbf{\te{ప్రశాంతమైన}}, \emph{praśāntamaina}.) A confluence? (\textbf{\te{సంగమం}}, \emph{saṅgamaṁ}.) A dam? (\textbf{\te{ఆనకట్ట}}, \emph{ānakaṭṭa}.) A reservoir? (\textbf{\te{జలాశయం}}, \emph{jalāśayaṁ}.) Irrigation? (\textbf{\te{నీటిపారుదల}}, \emph{nīṭipārudala}.) A marsh? (\textbf{\te{చిత్తడి}}, \emph{cittaḍi}.) The climate? (\textbf{\te{శీతోష్ణస్థితి}}, \emph{śītōṣṇasthiti}.) Temperature? (\textbf{\te{ఉష్ణోగ్రత}}, \emph{uṣṇōgrata}.) Humidity? (\textbf{\te{తేమ}}, \emph{tēma}.) Muggy heat? (\textbf{\te{ఉక్కపోత}}, \emph{ukkapōta}.) To heat up? (\textbf{\te{వేడెక్కు}}, \emph{vēḍekku}.) The monsoon? (\textbf{\te{ఋతుపవనం}}, \emph{ṛtupavanaṁ}.) Rainfall? (\textbf{\te{వర్షపాతం}}, \emph{varṣapātaṁ}.) A spell of drizzle? (\textbf{\te{ముసురు}}, \emph{musuru}.) A shower of rain? (\textbf{\te{జల్లు}}, \emph{jallu}.) To get soaked? (\textbf{\te{తడిసిపోవు}}, \emph{taḍisipōvu}.)
\end{tcolorbox}
